% mode: ini
% Aimed at "Resume the page builder after an output routine has come to an end": a closing brace that was backed up (the number
% of \count0=1 is ended by it), so that a backed-up list at its end is popped first; "Output routine didn't use all of \box255"
% with box_error; held-over insertions, material the routine left on its list, and contributions, going back in order
% through page_tail; flush of page_disc; a paragraph left open in the routine.
\input prelude
\errorcontextlines=0
\showboxbreadth=100
\hsize=100pt \vsize=20pt
\def\mk#1{\hrule height 15pt depth 0pt width #1pt}
% The brace of the routine is backed up by the scan of a number.
\output={\setbox0\box255 \deadcycles=0 \count0=1}
\mk1 \penalty-10000 \showlists
\output={\setbox0\box255 \deadcycles=0 \dimen0=1pt}
\mk2 \penalty-10000 \showlists
\output={\setbox0\box255 \deadcycles=0 \ifnum0=1 \fi}
\mk3 \penalty-10000 \showlists
% Box 255 not emptied.  Twice, one for a non-void left in a group.
\output={\deadcycles=0 }
\mk4 \penalty-10000 \showlists
\output={\global\setbox255\vbox{\hrule}\deadcycles=0 }
\mk5 \penalty-10000 \showlists
\output={\setbox255\vbox{\hrule}\deadcycles=0 }
\mk6 \penalty-10000 \showlists
% What the routine leaves on its list comes back, before the contributions.
\output={\setbox0\box255 \deadcycles=0 \showlists\hrule height 1pt width 7pt \hbox{left}\showlists}
\mk7 \penalty-10000 \hrule height 2pt width 8pt \showlists
\output={\setbox0\box255 \deadcycles=0 }
\penalty-10000 \showlists
% Insertions held over by \holdinginserts: they stand on the page's list, and contributions follow.
\holdinginserts=1
\skip100=0pt \count100=1000 \dimen100=100pt
\output={\showbox255 \global\setbox2\box255 \deadcycles=0 \showlists}
\mk9 \insert100{\hrule height 1pt}\penalty-10000 \hrule height 3pt\showlists
\penalty-10000 \showlists
\holdinginserts=0
\output={\showbox255 \setbox0\box255 \showbox100 \global\setbox2\box100 \deadcycles=0 }
\penalty-10000 \showlists
\mk1 \insert100{\hrule height 5pt}\penalty-10000 \showlists
% An insertion that splits: the remainder is a contribution to the next page, by way of link(page_head).
\dimen100=3pt \vsize=30pt
\mk1 \insert100{\hrule height 2pt\vskip1pt\hrule height 2pt\vskip1pt\hrule height 2pt}\hrule height 5pt\penalty-10000 \showlists
\penalty-10000 \showlists
\dimen100=100pt \vsize=20pt
% A paragraph open at the routine's end, with \partokencontext at its default.
\output={\setbox0\box255 \deadcycles=0 \noindent x\showlists}
\mk2 \penalty-10000 \showlists
\count1=0
\output={\setbox0\box255 \deadcycles=0 \ifnum\count1<2 \global\advance\count1 by 1 \hrule\vskip1pt\fi}
\mk2 \penalty-10000 \showlists
\mk2 \penalty-10000 \showlists
% page_disc: discarded items saved, e-TeX's page_disc is flushed here when not used.
\output={\setbox0\box255 \deadcycles=0 }
\mk3 \vskip1pt\penalty-10000 \vskip2pt \kern3pt \showlists\hrule height 1pt \showlists
\penalty-10000
\showlists
\output={\setbox0\box255 \deadcycles=0 }
\end
